\documentclass[11pt,a4paper]{article}
\usepackage[utf8]{inputenc}
\usepackage[T1]{fontenc}
\usepackage[spanish]{babel}
\usepackage{amsmath,amssymb}
\usepackage{geometry}
\usepackage{booktabs}
\usepackage{hyperref}
\geometry{margin=2.3cm}
\hypersetup{colorlinks=true,linkcolor=black,urlcolor=blue}

\title{\texttt{SimpleAssist3D}: arquitectura reactiva de extrusión espacial}
\author{PerfectPatch --- Minecraft Forge 1.8.9}
\date{\today}

\begin{document}
\maketitle

\begin{abstract}
\texttt{SimpleAssist3D} es un ejecutor independiente de parches 3D. No reutiliza
el cursor terminal--origen nativo de \texttt{SimpleAssist} 1D: mantiene su propia
captura, selección reactiva, créditos, confirmación y watchdogs. Su objetivo es
transformar un ancla real de obsidiana/bedrock y un \texttt{PatchCrumb} en una
secuencia corta de placements físicamente verificables, sin modificar la cámara
local ni extender el alcance del servidor.
\end{abstract}

\tableofcontents

\section{Modelo geométrico}

Sea $\mathbf e$ la posición de los ojos, $\mathbf l$ el vector de mirada y
$\mathbf q$ el centro de la cara soporte elegida. Para una trayectoria Norte/Sur,
el crumb define el plano objetivo
\begin{equation}
  X=X_c,\qquad Y=Y_c,
\end{equation}
por lo que $Z$ representa profundidad. Para una trayectoria Este/Oeste, el plano
es $Z=Z_c, Y=Y_c$, con $X$ como profundidad. La distancia de progreso usada por el
cursor es
\begin{equation}
 D_g(\mathbf p)=
 \begin{cases}
 |p_x-X_c|+|p_y-Y_c|,& \text{N/S},\\
 |p_z-Z_c|+|p_y-Y_c|,& \text{E/O}.
 \end{cases}
\end{equation}
Un placement ordinario debe reducir $D_g$. Los placements de profundidad se
permiten únicamente como lead cinemático controlado (sección~\ref{sec:depth}).

\section{Captura y ciclo reactivo}

El módulo se ejecuta en \texttt{ClientTickEvent.Phase.START}. Para iniciar una
sesión necesita un crumb activo, un bloque aceptado en la mano y un ancla aceptada.
Con Asistencia espacial activa, la captura puede usar el cono expandido para
normalizar pequeños desajustes de puntería.

No se conserva una ruta A* larga. En cada tick se parte del último bloque
confirmado o del último bloque pendiente, se revalida el mundo actual y se escoge
el siguiente vecino. El estado persistente es deliberadamente pequeño:
\begin{center}
\begin{tabular}{ll}
\toprule
Estado & Propósito \\
\midrule
crumb, eje y revisión capturados & Meta geométrica vigente \\
último confirmado & Base estable del cursor \\
hasta dos pendientes & Ventana de extrusión y confirmación \\
ancla/paridad de profundidad & Conservación de la retícula agua--bloque \\
candidato retenido & Histéresis de selección \\
créditos y tiempos & Ritmo y recuperación de estancamiento \\
\bottomrule
\end{tabular}
\end{center}

\section{Cono espacial, reach y línea de visión}

La distancia longitudinal y la distancia perpendicular de una cara al rayo son
\begin{equation}
 \lambda=\frac{(\mathbf q-\mathbf e)\cdot\mathbf l}{\|\mathbf l\|^2},
 \qquad
 d_\perp=\left\|\mathbf q-\left(\mathbf e+\lambda\mathbf l\right)\right\|.
\end{equation}
La captura expandida requiere $\lambda\geq0$ y
\begin{equation}
 d_\perp\leq R_{\rm cone}=1.25.
\end{equation}
La normalización explícita por $\|\mathbf l\|^2$ protege el cálculo aunque el
vector de mirada no sea unitario.

La autorización física usa la posición de ojos viva, nunca la predicha:
\begin{equation}
 \|\mathbf q-\mathbf e\|\leq R_{\rm real},
 \qquad R_{\rm real}=\min(R_{\rm controller},4.5).
\end{equation}
Además, el soporte debe ser sólido y el raytrace debe llegar a ese soporte. Como
$\mathbf q$ cae exactamente sobre una frontera de AABB en 1.8.9, el extremo del
raytrace se introduce $\varepsilon=0.001$ bloques en el bloque soporte:
\begin{equation}
 \mathbf q_{\rm trace}=\mathbf q-\varepsilon\mathbf n_{\rm face}.
\end{equation}
Esto evita falsos misses en ángulos rasantes.

\section{Restricciones físicas duras}

Ningún score puede compensar una regla dura. Antes de emitir un placement se exige:
\begin{enumerate}
 \item El destino es reemplazable.
 \item Existe una de las seis caras soporte válida; en una ventana encadenada, el
       primer bloque pendiente puede ser soporte virtual del segundo.
 \item El hit vector pasa reach real, cono expandido, raytrace y línea de visión.
 \item El bloque lleno a colocar no intersecta la caja del jugador.
\end{enumerate}

Para la última condición se prueba
\begin{equation}
 \operatorname{AABB}_{\rm player}\oplus(0.05,0,0.05)
 \;\cap\;\operatorname{AABB}_{\rm destination}=\varnothing.
\end{equation}
La expansión es solo horizontal: no se expande $Y$, para no cancelar placements
legítimos bajo los pies.

Las rutas mixtas permanecen en $Y_c\pm1$. Una columna ya alineada con el plano
objetivo puede usar $Y_c\pm3$. En el cursor ordinario, la posición lateral del
jugador respecto al plano del crumb restringe la cara horizontal admisible; los
pasos verticales no heredan dicha restricción.

\section{Selección cinemática local}

El selector evalúa como máximo seis direcciones, por lo que su complejidad es
$O(1)$ por tick. La predicción solo ordena candidatos. Definiendo
$s^2=v_x^2+v_z^2$, el ojo predicho es
\begin{equation}
 \widehat{\mathbf e}=\mathbf e+f\mathbf v,\qquad
 f=\begin{cases}
 1,&s^2>0.015,\\
 0.5,&0.005<s^2\leq0.015,\\
 0,&s^2\leq0.005.
 \end{cases}
\end{equation}

Para una cara $c$, el coste implementado es
\begin{equation}
 \begin{split}
 C(c)={}&40\,D_g(c)+25\,D_{\rm depth}(c)
       +1.5\,\|\mathbf q_c-\widehat{\mathbf e}\|^2\\
      &+3\,\Delta\theta(c)+P_{\rm behind}(c),
 \end{split}
\end{equation}
donde
\begin{equation}
 \Delta\theta(c)=\arccos\!\left(
 \frac{(\mathbf q_c-\mathbf e)\cdot\mathbf l}
 {\|\mathbf q_c-\mathbf e\|\,\|\mathbf l\|}\right)
\end{equation}
y la penalización de inercia es
\begin{equation}
 P_{\rm behind}(c)=
 \begin{cases}
 -5\,\widehat{\mathbf v}_{xz}\cdot\widehat{(\mathbf q_c-\mathbf p)}_{xz},
 &\widehat{\mathbf v}_{xz}\cdot\widehat{(\mathbf q_c-\mathbf p)}_{xz}<0,\\
 0,&\text{en otro caso.}
 \end{cases}
\end{equation}
Los pesos, por tanto, son $w_g=40$, $w_d=25$, $w_r=1.5$,
$w_\theta=3$ y $w_b=5$.

\section{Lead de profundidad y separación}\label{sec:depth}

La profundidad no altera el plano final, pero no puede destruir la retícula del
parche. Todos los nodos aceptados respetan
\begin{equation}
 |d(\mathbf p)-d(\mathbf a)|\bmod2=0,
\end{equation}
donde $d$ es $Z$ para N/S y $X$ para E/O, y $\mathbf a$ es el ancla de sesión.
Un lead puede avanzar exactamente dos bloques en dirección de la velocidad de
profundidad. Solo se favorece si $|v_d|>0.005$ y nunca se aleja más de dos celdas
de $\mathbf a$. La celda intermedia queda como agua: no se coloca ni se inventa
como soporte virtual.

\section{Histéresis y degradación suave}

Un candidato retenido puede vivir dos ticks. En cada reutilización se recalculan
soporte, reach, línea de visión, hitbox y nodo guía contra el mundo vivo. Si falla
una regla dura, se descarta inmediatamente. Si sigue válido, un challenger lo
reemplaza solo cuando
\begin{equation}
 C_{\rm new}<0.85\,C_{\rm active}.
\end{equation}
Es decir, necesita al menos un $15\%$ de mejora. Cambios pequeños de ángulo o
distancia dentro del reach no purgan el candidato, evitando oscilación y jitter.

\section{Extrusión, créditos y red}

El bucket local se actualiza con CPS configurables ($1\ldots20$):
\begin{equation}
 K_t=\min\left(2,K_{t-1}+\Delta t\frac{\mathrm{CPS}}{20}\right),
 \qquad \text{emitir si }K_t\geq1,
\end{equation}
restando un crédito por placement. El motor espacial mantiene un segundo bucket
independiente, recargado en cada pre-motion y también limitado a dos interacciones.

Si existen créditos y soporte válidos, se permite una ventana de dos placements
pendientes. La secuencia se conserva en orden de red:
\begin{center}
\texttt{C05 Look} $\rightarrow$ \texttt{C08 Place} $\rightarrow$
\texttt{C05 Look} $\rightarrow$ \texttt{C08 Place}.
\end{center}
El C05 lleva yaw/pitch calculados hacia $\mathbf q$, pero no escribe
\texttt{rotationYaw} ni \texttt{rotationPitch} locales. El segundo C08 puede usar
el primer bloque aún pendiente como soporte únicamente si la línea hacia ese
soporte virtual está limpia.

\section{Confirmación, revisión y watchdogs}

Cada destino pendiente almacena block-id y tick de envío. Solo al observar ese
block-id en el mundo cliente se promueve a \emph{último confirmado}. Un pendiente
expira después de $80$ ticks; el cursor también se reinicia tras $80$ ticks sin
progreso o tras tres misses de solicitud. Durante \texttt{isUsingItem()}, el
watchdog de progreso se suspende, aunque un pendiente individual sigue expirando.

Al capturar se guarda la revisión de \texttt{PatchCrumbs}. Cuando cambia, se lee
el crumb nuevo. La sesión se conserva exclusivamente si eje, $Y$ y coordenada de
plano ($X$ para N/S, $Z$ para E/O) siguen iguales. En tal caso se adopta la
representación fresca sin perder pendientes ni créditos. Si el crumb expira o el
plano cambia, la sesión se reinicia antes de emitir otro paquete.

\section{Límites deliberados}

El módulo no mueve al jugador, no cambia de slot, no modifica la cámara local y
no aumenta el reach de Vanilla. Predicción, histéresis, lead y guía cambian qué
candidato es preferible; jamás convierten en válido un placement que falle una
restricción física.

\end{document}
